\chapter{Resultaten}

Voor de evaluatie van de proof-of-concept (POC) werd nadien een vragenlijst afgenomen bij vijf respondenten. De resultaten worden hieronder besproken op basis van ervaring, leerdoelen, functionaliteiten, begrip van hashing en veiligheid \& bewustwording.
Van de 5 deelnemers was er één persoon die helemaal nog geen ervaring had met een wachtwoordmanager (en heel weinig computerervaring), de anderen wel. Geen van de deelnemers had al een wachtwoordmanager bewust gedownload en geïnstalleerd maar gebruikte wel al een ingebouwde wachtwoordmanager, zoals bv de app "Wachtwoorden" op Iphone.

\section{Ervaring \& gebruiksvriendelijkheid}

De simulatie werd over het algemeen als gebruiksvriendelijk ervaren. Op een schaal van 1 tot 5 scoorde de begrijpelijkheid van de simulatie gemiddeld 4,4, wat aangeeft dat de meeste gebruikers de applicatie relatief intuïtief konden gebruiken.

De duidelijkheid van de instructies kreeg een iets lagere gemiddelde score van 4,0. Eén respondent gaf expliciet aan dat een overlay of pop-up instructies de ervaring zou verbeteren, wat wijst op mogelijke optimalisatie van de begeleiding binnen de interface. Het was handiger geweest om telkens de instructies centraal op het scherm te laten verschijnen ipv enkel links waardoor er meer afleiding was en er soms te snel geklikt werd op andere knoppen.

De simulatie werd door alle deelnemers unaniem als realistisch beoordeeld, met een score van 5/5. Dit suggereert dat de POC erin slaagt een geloofwaardige nabootsing te bieden van een echte wachtwoordmanager.

\section{Bereiken van leerdoelen}

Wat betreft het begrijpen van het doel van een wachtwoordmanager gaven drie van de vier respondenten aan dit volledig te begrijpen, terwijl één respondent slechts gedeeltelijk begrip rapporteerde.

Alle respondenten gaven aan dat ze (al dan niet met enige begeleiding) bereid zouden zijn om een echte wachtwoordmanager te gebruiken na het doorlopen van de simulatie. Dit wijst op een positieve impact op de adoptiebereidheid.

Bij open vragen rond kennisverwerving bleek dat de meeste respondenten het verschil tussen sterke en zwakke wachtwoorden correct konden beschrijven. Sterke wachtwoorden werden geassocieerd met lengte, complexiteit en variatie in tekens, terwijl zwakke wachtwoorden gekenmerkt werden door eenvoud, hergebruik en persoonlijke informatie.

Het concept van een masterwachtwoord bleek minder goed begrepen: slechts een minderheid kon dit correct uitleggen, wat aangeeft dat dit onderdeel extra aandacht vereist binnen de simulatie. Eén persoon kon aangeven dat dit het wachtwoord is om in de kluis te geraken en het enige dat je als gebruiker moet onthouden.

\section{Functionaliteiten van de wachtwoordmanager}

De functionaliteiten van de simulatie werden zeer positief geëvalueerd. De automatische invulfunctie (autofill) kreeg een gemiddelde score van 5/5, wat het nut en de gebruiksvriendelijkheid van deze feature benadrukt. 

Daarnaast gaf het merendeel van de respondenten aan te begrijpen hoe wachtwoorden worden opgeslagen in de kluis, hoewel één respondent dit niet begreep.

De wachtwoordgenerator werd eveneens als duidelijk en nuttig ervaren, met een gemiddelde score van 4,75/5. Ook hier werd aangegeven dat het gemmakkelijk en tijdsbesparend is. Hier gaf één persoon aan dat hij hierdoor de feeling wat verliest met zijn wachtwoorden maar begreep ook dat het nuttig is om zonder nadenken sterke wachtwoorden aan te maken.

\section{Begrip van hashing}

Het begrip van hashing na de simulatie was gemengd. Slechts één respondent gaf aan dit volledig te begrijpen, terwijl anderen dit gedeeltelijk of niet begrepen.

Desondanks konden de meeste respondenten wel correct aangeven waarom wachtwoorden gehasht worden opgeslagen, namelijk om veiligheidsredenen en om het oorspronkelijke wachtwoord niet rechtstreeks te bewaren.

Bij het vergelijken van hashes merkten alle respondenten correct op dat kleine wijzigingen (één cijfer/letter/teken) in het wachtwoord leiden tot volledig verschillende hashes, wat een belangrijk leerdoel van de simulatie bevestigt.

\section{Veiligheid en bewustwording}

De resultaten tonen een variatie in huidig wachtwoordgebruik: sommige respondenten gebruiken reeds unieke wachtwoorden, terwijl anderen nog steeds (gedeeltelijk) hetzelfde wachtwoord hergebruiken.

Na de simulatie gaven de meeste respondenten aan zich bewuster te zijn van de risico’s van zwakke wachtwoorden, met scores tot 5/5. Eén respondent gaf aan reeds vooraf een hoog bewustzijn te hebben.

Daarnaast zouden bijna alle respondenten een wachtwoordmanager aanbevelen aan anderen, wat wijst op een positieve perceptie van het nut van dergelijke tools.

\section{Verbeterpunten}

Hoewel de algemene feedback positief was, werden enkele verbeterpunten geïdentificeerd:

\begin{itemize}
\item Instructies kunnen duidelijker, bijvoorbeeld via een pop-up of overlay (beperken van te vroeg op knoppen te klikken)
\item Visuele elementen zoals layout en pijltjes kunnen worden geoptimaliseerd (grotere pijl, andere kleur voor pijl)
\item Navigatie (zoals een “vorige” knop) kan verbeterd worden (werking nog niet geoptimaliseerd)
\item Nog meer inzetten op risico's van zwak wachtwoordgebruik zodat beginnende gebruikers nog meer gemotiveerd worden om een wachtwoordmanager te gebruiken
\item Info geven over welke wachtwoordmanagers er bestaan en of deze betalend of gratis zijn.
\end{itemize}

De duur van de tutorial werd door de meeste respondenten geschat op ongeveer 10 minuten, en de moeilijkheidsgraad werd doorgaans als “juist goed” ervaren, met één uitzondering waarbij extra begeleiding gewenst was.